% ============================================================================
%  第一章 绪论（教学诊断版示例）
%  包含 6 处典型错误标注，对应 student-pitfalls.md 的 P-1.1 ~ P-1.6
% ============================================================================

\chapter{绪论}
\label{chap:intro}

% ----------------------------------------------------------------------------
\section{研究背景与意义}
\label{sec:background}

\par 随着工业互联网（Industrial Internet of Things, IIoT）与人工智能技术的快
速发展，SCADA（Supervisory Control and Data Acquisition，监控与数据采集）系
统作为国家关键基础设施的神经中枢，正面临日益严峻的网络安全挑战
\cite{ref:scada_survey_2022,ref:iiot_security_2023}。根据国家工业信息安全
发展研究中心发布的报告，2024 年全球工业领域网络攻击事件同比增长 27\%，其中
针对能源、水务等 SCADA 系统的定向攻击占比超过 35\%。

\pitfall{P-1.1 数据无出处}{学生常引"据某报告"、"相关数据显示"等模糊表述，无
可追溯的数据来源；老师评阅时一票否决"可信度"项。建议：所有数据必须有 [n]
标注，引用具体机构、年份、报告名。}

\par SCADA 系统广泛应用于燃气管道、电力调度、轨道交通、自来水处理等关键领域，
其安全性直接关系到国计民生与公共安全。2010 年伊朗"震网"病毒、2015 年乌克兰
电网攻击、2021 年美国 Colonial Pipeline 输油管道勒索事件，无一不揭示了 SCADA
系统一旦失守可能造成的灾难性后果。

\good{注意本段：使用"震网"、"乌克兰电网"、"Colonial Pipeline"三个递进案例，
从国家→地区→企业层面论证，逻辑递进清晰。}

\par 在此背景下，入侵检测（Intrusion Detection System, IDS）作为 SCADA 安全的
最后一道防线，已成为学术界与工业界共同关注的研究热点。然而，SCADA 工控场景
对 IDS 提出了独特而严苛的要求：

\begin{itemize}
  \item \textbf{实时性}：工业控制周期通常在毫秒级，要求检测延迟 $<$ 10\,ms；
  \item \textbf{高准确率}：误报会导致运维人员"狼来了"疲劳，漏报则可能造成实
        际安全事故；
  \item \textbf{可部署性}：SCADA 现场设备（PLC、RTU）算力与存储有限，模型
        必须轻量化。
\end{itemize}

% ----------------------------------------------------------------------------
\section{国内外研究现状}
\label{sec:related}

\par 现有 SCADA 入侵检测研究大致可分为三类：基于规则与签名的方法、基于传统
机器学习的方法、基于深度学习的方法。

\subsection{基于规则与签名的方法}
\label{sec:rule_based}

\par 早期研究多采用专家系统与特征库匹配 \cite{ref:snort_1999}。该方法对已知
攻击检测准确率高、可解释性强，但完全依赖人工维护规则库，对未知攻击
（zero-day）几乎无能为力。

\subsection{基于传统机器学习的方法}
\label{sec:ml_based}

\par 随着 KDD Cup 99、NSL-KDD、UNSW-NB15 等入侵检测基准数据集的发布，基于
传统机器学习的方法成为研究主流。支持向量机（SVM）\cite{ref:svm_ids_2008}、
随机森林（Random Forest）\cite{ref:rf_ids_2010}、XGBoost
\cite{ref:xgboost_ids_2019} 等模型在多个数据集上取得了 F1 $\geq$ 0.85 的检测
性能。

\pitfall{P-1.2 文献堆砌}{学生常罗列 20+ 篇文献却不评价。"某某\cite{x}提出了
...，某某\cite{y}改进了..."——这种写法老师会扣"文献综合分析"分。建议：每段
文献综述末尾必须有 1 句评价，指出该方法的局限性，引出下一个小节。}

\subsection{基于深度学习的方法}
\label{sec:dl_based}

\par 近年来，深度学习在时序建模领域的突破为入侵检测带来了新思路。文献
\cite{ref:cnn_ids_2017} 首次将 1D-CNN 应用于网络流量分类，取得了比 SVM 高
3\% 的准确率；文献 \cite{ref:lstm_ids_2018} 提出 LSTM-IDS，利用长短期记忆
网络捕捉时序依赖；文献 \cite{ref:transformer_ids_2021} 进一步引入注意力
机制，在多个公开数据集上达到 SOTA。

然而，上述工作大多基于通用网络入侵数据集（KDD/NSL-KDD/UNSW-NB15），鲜有
针对 SCADA 协议语义（如 Modbus、DNP3）的专门优化。

\pitfall{P-1.3 缺乏协议层切入}{这是本领域最大空白。SCADA 与传统 IT 网络的根
本区别在协议层；学生若不强调"协议语义"，会和通用 IDS 论文混为一谈。老师会
问："你的工作和 NSL-KDD 上的工作有什么本质区别？"——这就是回答点。}

% ----------------------------------------------------------------------------
\section{现有研究不足}
\label{sec:gap}

\par 综合分析国内外研究现状，当前 SCADA 入侵检测领域存在以下不足：

\begin{enumerate}
  \item \textbf{特征体系缺乏协议语义}：现有研究多沿用通用网络入侵特征，未
        充分利用 Modbus 主-从交互结构等协议层先验；
  \item \textbf{深度学习模型未稳定超越树模型}：在多个公开数据集上，LightGBM
        /XGBoost 仍占据 SOTA，DL 的优势尚未充分释放；
  \item \textbf{模型部署可行性研究不足}：现有工作多报告 GPU 推理延迟，对
        MCU/PLC 等受限设备的部署研究较少。
\end{enumerate}

\good{本段是"研究意义"和"研究内容"之间的桥梁，三个"不足"恰好对应后文
三个"创新点"——这是 1:1 映射的最佳实践。}

% ----------------------------------------------------------------------------
\section{研究内容与创新点}
\label{sec:contribution}

\par 针对上述不足，本文以 IanArffDataset 数据集为研究对象，开展从数据特征、
模型架构到边缘部署的全栈研究，主要创新点如下：

\begin{enumerate}
  \item \textbf{基于 SCADA 协议语义的领域特征工程}（对应 \chapref{chap:feat}）。
        分析 Modbus 4 行主-从交互结构，提出 27 维协议语义特征，使所有模型
        Macro-F1 一致提升 0.2--0.7 个百分点；
  \item \textbf{TCN-SE 轻量入侵检测模型}（对应 \chapref{chap:tcn}）。在 27
        种时序模型的横向对比基础上，融合 SE 通道注意力，TCN-SE 在测试集
        取得 Macro-F1 = 0.8340，首次在该任务上超越 LightGBM；
  \item \textbf{面向 MCU 的极限压缩与混合精度量化}（对应 \chapref{chap:deploy}）。
        提出 4{,}237 参数的 TCN V19 变体与权重 INT8 / 激活 FP32 混合方案，
        在 Cortex-M4 平台上实现 5.5\,KB 模型、0.47\,ms 推理。
\end{enumerate}

\pitfall{P-1.4 创新点无量化数据}{"创新点"是论文最核心的部分，但学生常写成
"提出了一种新方法"——必须每个创新点都有可量化的指标提升。本模板已给出
+0.2--0.7pp / 0.8340 / 5.5KB / 0.47ms 等具体数字。}

\pitfall{P-1.5 创新点对应关系模糊}{"对应第 X 章"必须明确写出，且在第 X 章
开头必须有"本章对应创新点 N"的回指。老师会查"前后呼应"。}

% ----------------------------------------------------------------------------
\section{论文组织结构}
\label{sec:structure}

\par 本文章节安排如下：

\begin{itemize}
  \item \chapref{chap:intro} 绪论：介绍研究背景、国内外现状与本文创新点；
  \item \chapref{chap:related} 相关理论与技术基础：SCADA 入侵检测任务定义、
        TCN 与注意力机制原理、量化基础；
  \item \chapref{chap:feat} SCADA 协议特征工程：详细阐述 27 维特征的设计
        动机与对比实验；
  \item \chapref{chap:tcn} TCN-SE 入侵检测模型：模型架构、27 模型横评、消
        融实验；
  \item \chapref{chap:deploy} 极限压缩与混合精度量化：V19 变体设计、INT8
        量化方案、MCU 部署实测；
  \item \chapref{chap:conclusion} 总结与展望。
\end{itemize}

\good{每章一句话点明"做什么"，避免冗长。}
